\section{Problem Definition and Stakeholder Context}
\label{sec:problem}

Phase 1 left one finding standing above the rest: a late delivery raises the
probability of a 1--2 star review from 9.2\% to 54.0\%, a 5.9-fold increase
monotone in how late the parcel is. Phase 2 turns that sentence into a pair of
predictions placed at two points in an order's life.

\begin{table}[htbp]
\centering
\caption{The two problems. They share one feature pipeline, one model-family
budget and one dataset; what separates them is \emph{when} the model runs.}
\label{tab:problems}
\small
\begin{tabular}{@{}p{1.9cm}p{6.2cm}p{6.2cm}@{}}
\toprule
 & \textbf{Problem 1 --- delivery performance} & \textbf{Problem 2 --- satisfaction risk} \\
\midrule
Question & How long will this order take, and will it miss its promise? &
Will this delivered order draw a 1--2 star review? \\[2pt]
Moment & \code{at\_checkout} --- the order has just been placed &
\code{at\_delivery} --- the parcel has arrived, the review is not yet written \\[2pt]
Stakeholder & Operations and logistics & Customer experience \\[2pt]
Task & \textbf{1A regression}: lead time in days, on MAE.
\textbf{1B classification}: delivered after the promised date, on PR-AUC against
its no-skill floor &
\textbf{Classification}: \code{review\_score} $\le 2$, on PR-AUC against its
floor and balanced accuracy \\[2pt]
Action & Re-quote the promise, warn the buyer, expedite freight &
Reach out before the review lands \\
\bottomrule
\end{tabular}
\end{table}

The brief permits two problems or one framed twice; we take two, because they are
halves of one decision chain. Operations acts before the parcel ships; customer
experience acts after it arrives and before the customer writes anything.
Problem 1's output is, in substance, Problem 2's most important input.

Problem 1 is the brief's \emph{dual framing}: one underlying quantity, two
outputs, so the late-order flag \emph{is} the regression's answer compared
against the promised date. Five classifiers trained directly on the binary label
run alongside as a control arm, testing whether the flag needs a model of its
own; Section~\ref{sec:results-late} reports that comparison, including a
disagreement between the selection protocol and the forward window that is itself
one of this report's findings.

\paragraph{Targets.} All three come from the provided tables alone.
\emph{Lead time} is \code{order\_delivered\_customer\_date} $-$
\code{order\_purchase\_timestamp} in days, defined for the 97.0\% of orders that
reached the customer. \emph{Late} is delivery strictly after
\code{order\_estimated\_delivery\_date}, which carries no time component, so
lateness is accurate to about a day. \emph{Low review} is \code{review\_score}
$\le 2$, defined for the 99.2\% carrying a review. Undefined rows are \code{NaN}
and excluded, never imputed to a negative label: the event has not happened
rather than gone missing, and a silent customer is not a satisfied one.

\paragraph{Success criteria, registered in advance.}
Expected ranges for all six headline metrics were written into
\code{config.EXPECTED\_BANDS} \emph{before} any Phase 2 model was tuned, from the
Phase 1 results and untuned scoping probes (Appendix~\ref{app:bands}). They are
the targets the work is judged against and they are a leak alarm --- because here
the signature of leakage is not a subtle improvement but an implausible one, as
Phase 1 showed by adding one forbidden column to the late-delivery model and
watching ROC-AUC go from 0.589 to 1.000. A result landing far \emph{above} its
band is investigated before it is reported. All six landed inside; none above.

The ceiling deserves stating as plainly as the target: lateness here is volatile
rather than trending --- 1.4\% to 21.4\% month to month with Black Friday
congestion and the May 2018 truckers' strike --- and no order-level feature
available at checkout can anticipate a national road blockade. A late-delivery
classifier reporting 95\% accuracy here would be the majority-class rule or a
leak.
